\IfFileExists{stacks-project.cls}{%
\documentclass{stacks-project}
}{%
\documentclass{amsart}
}
\usepackage{amssymb}

% For dealing with references we use the comment environment
\usepackage{verbatim}
\usepackage{fancyvrb}
\usepackage{pmboxdraw}
\newenvironment{reference}{\comment}{\endcomment}
%\newenvironment{reference}{}{}
\newenvironment{slogan}{\comment}{\endcomment}
\newenvironment{history}{\comment}{\endcomment}

% For source figures retained as assets
\usepackage{graphicx}

% For commutative diagrams we use Xy-pic
\usepackage[all]{xy}

% We use 2cell for 2-commutative diagrams.
\xyoption{2cell}
\UseAllTwocells

% We use multicol for the list of chapters between chapters
\usepackage{multicol}

% This is generally recommended for better output
\usepackage{lmodern}
\usepackage[T1]{fontenc}
\usepackage{textalpha}

% For cross-file-references
% We prefer the snapshot tmp/xref/#2.aux,
% which latexmkrc replaces atomically, so
% that documents compiled simultaneously
% never read a half-written .aux file.
% The last argument keeps links to #2.pdf.
\usepackage{xr-hyper}
\newcommand{\maybeexternaldocument}[2]{%
  \IfFileExists{#2.aux}{%
    \externaldocument[#1]{#2}%
  }{%
    \PackageWarning{stack}{Missing #2.aux; external references with prefix #1 unavailable}%
  }%
}
\let\externaldocumentorig\externaldocument
\renewcommand{\externaldocument}[2][]{%
  \IfFileExists{tmp/xref/#2.aux}{%
    \externaldocumentorig[#1]{tmp/xref/#2}[#2.pdf]%
  }{%
  \IfFileExists{#2.aux}{%
    \externaldocumentorig[#1]{#2}%
  }{%
    \PackageWarning{stack}{Missing #2.aux; external references with prefix #1 unavailable}%
  }%
  }%
}

% Package for hypertext links:
\usepackage{hyperref}

% For any local file, say "hello.tex" you want to link to please
% use \externaldocument[hello-]{hello}
\externaldocument[introduction-]{introduction}
\externaldocument[conventions-]{conventions}
\externaldocument[sets-]{sets}
\externaldocument[categories-]{categories}
\externaldocument[topology-]{topology}
\externaldocument[sheaves-]{sheaves}
\externaldocument[sites-]{sites}
\externaldocument[stacks-]{stacks}
\externaldocument[fields-]{fields}
\externaldocument[algebra-]{algebra}
\externaldocument[brauer-]{brauer}
\externaldocument[homology-]{homology}
\externaldocument[derived-]{derived}
\externaldocument[simplicial-]{simplicial}
\externaldocument[more-algebra-]{more-algebra}
\externaldocument[smoothing-]{smoothing}
\externaldocument[modules-]{modules}
\externaldocument[sites-modules-]{sites-modules}
\externaldocument[injectives-]{injectives}
\externaldocument[cohomology-]{cohomology}
\externaldocument[sites-cohomology-]{sites-cohomology}
\externaldocument[dga-]{dga}
\externaldocument[dpa-]{dpa}
\externaldocument[sdga-]{sdga}
\externaldocument[hypercovering-]{hypercovering}
\externaldocument[schemes-]{schemes}
\externaldocument[constructions-]{constructions}
\externaldocument[properties-]{properties}
\externaldocument[morphisms-]{morphisms}
\externaldocument[coherent-]{coherent}
\externaldocument[divisors-]{divisors}
\externaldocument[limits-]{limits}
\externaldocument[varieties-]{varieties}
\externaldocument[topologies-]{topologies}
\externaldocument[descent-]{descent}
\externaldocument[perfect-]{perfect}
\externaldocument[more-morphisms-]{more-morphisms}
\externaldocument[flat-]{flat}
\externaldocument[groupoids-]{groupoids}
\externaldocument[more-groupoids-]{more-groupoids}
\externaldocument[etale-]{etale}
\externaldocument[chow-]{chow}
\externaldocument[intersection-]{intersection}
\externaldocument[pic-]{pic}
\externaldocument[weil-]{weil}
\externaldocument[adequate-]{adequate}
\externaldocument[dualizing-]{dualizing}
\externaldocument[duality-]{duality}
\externaldocument[discriminant-]{discriminant}
\externaldocument[derham-]{derham}
\externaldocument[local-cohomology-]{local-cohomology}
\externaldocument[algebraization-]{algebraization}
\externaldocument[curves-]{curves}
\externaldocument[resolve-]{resolve}
\externaldocument[models-]{models}
\externaldocument[functors-]{functors}
\externaldocument[equiv-]{equiv}
\externaldocument[pione-]{pione}
\externaldocument[etale-cohomology-]{etale-cohomology}
\externaldocument[proetale-]{proetale}
\externaldocument[relative-cycles-]{relative-cycles}
\externaldocument[more-etale-]{more-etale}
\externaldocument[trace-]{trace}
\externaldocument[crystalline-]{crystalline}
\externaldocument[spaces-]{spaces}
\externaldocument[spaces-properties-]{spaces-properties}
\externaldocument[spaces-morphisms-]{spaces-morphisms}
\externaldocument[decent-spaces-]{decent-spaces}
\externaldocument[spaces-cohomology-]{spaces-cohomology}
\externaldocument[spaces-limits-]{spaces-limits}
\externaldocument[spaces-divisors-]{spaces-divisors}
\externaldocument[spaces-over-fields-]{spaces-over-fields}
\externaldocument[spaces-topologies-]{spaces-topologies}
\externaldocument[spaces-descent-]{spaces-descent}
\externaldocument[spaces-perfect-]{spaces-perfect}
\externaldocument[spaces-more-morphisms-]{spaces-more-morphisms}
\externaldocument[spaces-flat-]{spaces-flat}
\externaldocument[spaces-groupoids-]{spaces-groupoids}
\externaldocument[spaces-more-groupoids-]{spaces-more-groupoids}
\externaldocument[bootstrap-]{bootstrap}
\externaldocument[spaces-pushouts-]{spaces-pushouts}
\externaldocument[spaces-chow-]{spaces-chow}
\externaldocument[groupoids-quotients-]{groupoids-quotients}
\externaldocument[spaces-more-cohomology-]{spaces-more-cohomology}
\externaldocument[spaces-simplicial-]{spaces-simplicial}
\externaldocument[spaces-duality-]{spaces-duality}
\externaldocument[formal-spaces-]{formal-spaces}
\externaldocument[restricted-]{restricted}
\externaldocument[spaces-resolve-]{spaces-resolve}
\externaldocument[formal-defos-]{formal-defos}
\externaldocument[defos-]{defos}
\externaldocument[cotangent-]{cotangent}
\externaldocument[examples-defos-]{examples-defos}
\externaldocument[algebraic-]{algebraic}
\externaldocument[examples-stacks-]{examples-stacks}
\externaldocument[stacks-sheaves-]{stacks-sheaves}
\externaldocument[criteria-]{criteria}
\externaldocument[artin-]{artin}
\externaldocument[quot-]{quot}
\externaldocument[stacks-properties-]{stacks-properties}
\externaldocument[stacks-morphisms-]{stacks-morphisms}
\externaldocument[stacks-limits-]{stacks-limits}
\externaldocument[stacks-cohomology-]{stacks-cohomology}
\externaldocument[stacks-perfect-]{stacks-perfect}
\externaldocument[stacks-introduction-]{stacks-introduction}
\externaldocument[stacks-more-morphisms-]{stacks-more-morphisms}
\externaldocument[stacks-geometry-]{stacks-geometry}
\externaldocument[moduli-]{moduli}
\externaldocument[moduli-curves-]{moduli-curves}
\externaldocument[examples-]{examples}
\externaldocument[exercises-]{exercises}
\externaldocument[guide-]{guide}
\externaldocument[desirables-]{desirables}
\externaldocument[coding-]{coding}
\externaldocument[obsolete-]{obsolete}
\externaldocument[fdl-]{fdl}
\externaldocument[index-]{index}

% Heap Project documents
\externaldocument[k3-]{k3}
\externaldocument[torelli-]{torelli}
\externaldocument[linear-algebra-]{linear-algebra}
\externaldocument[groups-]{groups}
\externaldocument[complex-analysis-]{complex-analysis}
\externaldocument[functional-analysis-]{functional-analysis}
\externaldocument[categories-]{categories}
\externaldocument[commutative-algebra-]{commutative-algebra}
\externaldocument[ringed-spaces-]{ringed-spaces}
\externaldocument[homological-algebra-]{homological-algebra}
\externaldocument[lie-algebras-]{lie-algebras}
\externaldocument[representation-theory-]{representation-theory}
\externaldocument[smooth-manifolds-]{smooth-manifolds}
\externaldocument[differential-topology-]{differential-topology}
\externaldocument[algebraic-topology-]{algebraic-topology}
\externaldocument[differential-geometry-]{differential-geometry}
\externaldocument[algebraic-geometry-]{algebraic-geometry}
\externaldocument[symplectic-geometry-]{symplectic-geometry}
\externaldocument[complex-geometry-]{complex-geometry}
\externaldocument[hodge-theory-]{hodge-theory}
\externaldocument[%
variations-of-hodge-structures-]{%
variations-of-hodge-structures}
\externaldocument[%
homological-mirror-symmetry-]{%
homological-mirror-symmetry}
\maybeexternaldocument{deformations-}{deformations}
\externaldocument[vertex-algebras-]{vertex-algebras}
\externaldocument[springer-]{springer}
\externaldocument[infty-categories-]{infty-categories}
\externaldocument[seminars-differential-geometry-]{%
seminars-differential-geometry}
\externaldocument[seminars-complex-geometry-]{%
seminars-complex-geometry}
\externaldocument[seminars-algebraic-geometry-]{%
seminars-algebraic-geometry}
\externaldocument[seminars-analysis-pde-]{%
seminars-analysis-pde}
\externaldocument[surfaces-]{surfaces}
\externaldocument[birational-maps-commutative-algebra-]{birational-maps-commutative-algebra}
\externaldocument[cremona-transformations-]{cremona-transformations}
\externaldocument[derived-categories-of-sheaves-]{derived-categories-of-sheaves}
\externaldocument[geometric-invariant-theory-]{geometric-invariant-theory}
\externaldocument[geometric-stability-conditions-and-group-actions-]{geometric-stability-conditions-and-group-actions}
\externaldocument[moduli-spaces-of-sheaves-]{moduli-spaces-of-sheaves}
\externaldocument[bridgeland-stability-general-theory-]{bridgeland-stability-general-theory}
\externaldocument[hilbert-schemes-]{hilbert-schemes}
\externaldocument[noncommutative-abelian-surfaces-and-kummer-type-hyperkahler-varieties-]{noncommutative-abelian-surfaces-and-kummer-type-hyperkahler-varieties}
\externaldocument[mumford-tate-groups-in-hodge-theory-]{mumford-tate-groups-in-hodge-theory}
\externaldocument[hodge-birational-atoms-2026-]{hodge-birational-atoms-2026}
\externaldocument[xxii-egd-2026-]{%
xxii-egd-2026}
\externaldocument[pde-]{pde}
\externaldocument[probability-]{probability}

\externaldocument[linguistics-]{linguistics}

\externaldocument[%
16th-alga-meeting-2026-]{%
16th-alga-meeting-2026}

\externaldocument[%
iv-jornadas-lefschetz-2026-]{%
iv-jornadas-lefschetz-2026}

\externaldocument[reading-the-masters-]{%
reading-the-masters}

\externaldocument[real-analysis-]{%
real-analysis}

\externaldocument[optimal-transport-]{%
optimal-transport}

\externaldocument[history-]{history}

\externaldocument[seiberg-witten-]{seiberg-witten}

\externaldocument[history-philosophy-biology-]{history-philosophy-biology}

\externaldocument[classical-mechanics-]{classical-mechanics}
\externaldocument[electromagnetism-]{electromagnetism}
\externaldocument[waves-and-optics-]{waves-and-optics}
\externaldocument[wif-]{waves-in-fluids-iv}
\externaldocument[thermal-physics-]{thermal-physics}
\externaldocument[quantum-mechanics-]{quantum-mechanics}
\externaldocument[special-relativity-]{special-relativity}
\externaldocument[physics-]{physics}

\externaldocument[chemistry-]{chemistry}

\externaldocument[molecular-cell-biology-]{molecular-cell-biology}

\externaldocument[microbiology-]{microbiology}

\externaldocument[protistology-]{protistology}

\externaldocument[mycology-]{mycology}

\externaldocument[botany-]{botany}

\externaldocument[zoology-]{zoology}

\externaldocument[animal-histology-]{animal-histology}

\externaldocument[animal-physiology-]{animal-physiology}

\externaldocument[systematics-biogeography-]{systematics-biogeography}

\externaldocument[ecology-natural-resources-]{ecology-natural-resources}

\externaldocument[earth-sciences-]{earth-sciences}

\externaldocument[biostatistics-]{biostatistics}

\externaldocument[mathematical-systems-biology-]{mathematical-systems-biology}

\externaldocument[cimpa-el-salvador-2024-]{cimpa-el-salvador-2024}

\externaldocument[anthropology-]{anthropology}

\externaldocument[yang-mills-theory-]{yang-mills-theory}

\externaldocument[seminars-algebra-]{seminars-algebra}

\externaldocument[seminars-physics-]{seminars-physics}

\externaldocument[seminars-computing-]{seminars-computing}

\externaldocument[oktobermat-2026-]{oktobermat-2026}

\externaldocument[saet-2026-]{saet-2026}

\externaldocument[seminars-symplectic-geometry-]{seminars-symplectic-geometry}

\externaldocument[seminars-dynamics-]{seminars-dynamics}

\externaldocument[oktobermat-2024-]{oktobermat-2024}

\externaldocument[quantum-groups-]{quantum-groups}

\externaldocument[legal-2027-]{legal-2027}

\externaldocument[derived-categories-]{derived-categories}

% Theorem environments.
%
\theoremstyle{plain}
\newtheorem{theorem}[subsection]{Theorem}
\newtheorem{proposition}[subsection]{Proposition}
\newtheorem{lemma}[subsection]{Lemma}
\newtheorem{corollary}[subsection]{Corollary}

\theoremstyle{definition}
\newtheorem{definition}[subsection]{Definition}
\newtheorem{example}[subsection]{Example}
\newtheorem{exercise}[subsection]{Exercise}
\newtheorem{situation}[subsection]{Situation}

\theoremstyle{remark}
\newtheorem{remark}[subsection]{Remark}
\newtheorem{remarks}[subsection]{Remarks}

\numberwithin{equation}{subsection}

% Macros
%
\def\lim{\mathop{\mathrm{lim}}\nolimits}
\def\colim{\mathop{\mathrm{colim}}\nolimits}
\def\Spec{\mathop{\mathrm{Spec}}}
\def\Hom{\mathop{\mathrm{Hom}}\nolimits}
\def\Ext{\mathop{\mathrm{Ext}}\nolimits}
\def\SheafHom{\mathop{\mathcal{H}\!\mathit{om}}\nolimits}
\def\SheafExt{\mathop{\mathcal{E}\!\mathit{xt}}\nolimits}
\def\Sch{\mathit{Sch}}
\def\Mor{\mathop{\mathrm{Mor}}\nolimits}
\def\Ob{\mathop{\mathrm{Ob}}\nolimits}
\def\Sh{\mathop{\mathit{Sh}}\nolimits}
\def\NL{\mathop{N\!L}\nolimits}
\def\CH{\mathop{\mathrm{CH}}\nolimits}
\def\proetale{{pro\text{-}\acute{e}tale}}
\def\etale{{\acute{e}tale}}
\def\QCoh{\mathit{QCoh}}
\def\Ker{\mathop{\mathrm{Ker}}}
\def\Im{\mathop{\mathrm{Im}}}
\def\Coker{\mathop{\mathrm{Coker}}}
\def\Coim{\mathop{\mathrm{Coim}}}

% Boxtimes
%
\DeclareMathSymbol{\boxtimes}{\mathbin}{AMSa}{"02}

%
% Macros for moduli stacks/spaces
%
\def\QCohstack{\mathcal{QC}\!\mathit{oh}}
\def\Cohstack{\mathcal{C}\!\mathit{oh}}
\def\Spacesstack{\mathcal{S}\!\mathit{paces}}
\def\Quotfunctor{\mathrm{Quot}}
\def\Hilbfunctor{\mathrm{Hilb}}
\def\Curvesstack{\mathcal{C}\!\mathit{urves}}
\def\Polarizedstack{\mathcal{P}\!\mathit{olarized}}
\def\Complexesstack{\mathcal{C}\!\mathit{omplexes}}
% \Pic is the operator that assigns to X its picard group, usage \Pic(X)
% \Picardstack_{X/B} denotes the Picard stack of X over B
% \Picardfunctor_{X/B} denotes the Picard functor of X over B
\def\Pic{\mathop{\mathrm{Pic}}\nolimits}
\def\Picardstack{\mathcal{P}\!\mathit{ic}}
\def\Picardfunctor{\mathrm{Pic}}
\def\Deformationcategory{\mathcal{D}\!\mathit{ef}}


\begin{document}

\title{LEGAL 2027}
\maketitle
\phantomsection
\label{section-phantom}
\tableofcontents

\medskip\noindent
{\bf Conference information.}

\href{%
https://sites.google.com/view/%
geoalgcompluff/legal%
}{LEGAL 2027 (Liberdade em Geometria
Alg\'ebrica)}

March 1--6, 2027.

\href{https://sitioassuncao.wordpress.com}{%
S\'itio Assun\c{c}\~ao},
Teresópolis, Rio de Janeiro, Brazil.

Information accessed November 20, 2026.
The official site was last updated on
November 19, 2026.

The official site lists no day, time or
room for any item as of November 20, 2026.
The sections follow the order of the
official abstracts page: minicourse,
invited speakers, special sessions
(mentors), contributed talks and posters.

\medskip\noindent
{\bf Organizers.}

Alex Abreu, Juliana Coelho, Marco Pacini
and Valeriano Lanza (UFF, Niter\'oi).

\section{From the Minimal Model Program to
Fano Conic Bundles}
\label{section-from-the-minimal-model}
%% speaker: Eleonora Anna Romano, Università
%% speaker: di Genova
%% place: LEGAL 2027, Sítio Assunção,
%% place: Teresópolis, RJ, Brazil

\medskip\noindent
{\bf Official presentation type.}

Minicourse.

\medskip\noindent
{\bf Speaker.}

Eleonora Anna Romano, Università di Genova.

\medskip\noindent
{\bf Schedule.}

Not listed on the official site as of
November 20, 2026.

\medskip\noindent
{\bf Official listing.}

\href{%
https://sites.google.com/view/%
geoalgcompluff/legal/%
legal-2027-resumos#h.c1v018hrntqh%
}{Programme entry.}

\medskip\noindent
{\bf Abstract.}

This minicourse is an introduction to
birational geometry and the Minimal Model
Program (MMP), with a focus on Fano
varieties. After an overview of the aims of
birational geometry and of the MMP, we
illustrate the theory in the case of
surfaces, before turning to Fano varieties in
arbitrary dimension and discussing their
basic properties and low-dimensional
examples. The last part of the course focuses
on conic bundles, one of the main tools for
studying and classifying Fano varieties
through Mori theory; we present structural
results for Fano conic bundles and their
applications to classification problems,
based on the speaker's own research.

Lecture 1: An introduction to the Minimal
Model Program. The classification problem up
to birational equivalence; a panoramic
overview of the Minimal Model Program and its
goals; the role of Fano varieties and conic
bundles within this classification program.

Lecture 2: The case of surfaces and the case
of Fano varieties. Mori theory for surfaces;
transition to Fano varieties in arbitrary
dimension and their basic properties, with
low-dimensional examples.

Lecture 3: Fano Conic bundles, Part I. Conic
bundles over Fano varieties, examples.
Geometric properties, discriminant locus,
factorization of non-elementary conic bundles
via birational geometry.

Lecture 4: Fano Conic bundles, Part II. A
characterization of some Fano 4-folds via
conic bundles, applications to
classification, open problems based on the
speaker’s own work.

\section{Title to be announced: Lyalya
Guseva}
\label{section-tba-guseva}
%% speaker: Lyalya Guseva, Université
%% speaker: Bourgogne Franche-Comté; HSE
%% speaker: University
%% place: LEGAL 2027, Sítio Assunção,
%% place: Teresópolis, RJ, Brazil

\medskip\noindent
{\bf Official presentation type.}

Invited speaker.

\medskip\noindent
{\bf Title.}

The official site lists the title as
``TBA'' as of November 20, 2026.

\medskip\noindent
{\bf Speaker.}

Lyalya Guseva, Université Bourgogne
Franche-Comté; HSE University.

\medskip\noindent
{\bf Schedule.}

Not listed on the official site as of
November 20, 2026.

\medskip\noindent
{\bf Official listing.}

\href{%
https://sites.google.com/view/%
geoalgcompluff/legal/%
legal-2027-resumos#h.59gubrdmgdxn%
}{Programme entry.}

\medskip\noindent
{\bf Abstract.}

No abstract is available on the
official site as of November 20, 2026.

\section{Title to be announced: Renato Vidal
Martins}
\label{section-tba-martins}
%% speaker: Renato Vidal Martins, UFMG
%% place: LEGAL 2027, Sítio Assunção,
%% place: Teresópolis, RJ, Brazil

\medskip\noindent
{\bf Official presentation type.}

Invited speaker.

\medskip\noindent
{\bf Title.}

The official site lists the title as
``TBA'' as of November 20, 2026.

\medskip\noindent
{\bf Speaker.}

Renato Vidal Martins, UFMG.

\medskip\noindent
{\bf Schedule.}

Not listed on the official site as of
November 20, 2026.

\medskip\noindent
{\bf Official listing.}

\href{%
https://sites.google.com/view/%
geoalgcompluff/legal/%
legal-2027-resumos#h.59gubrdmgdxn%
}{Programme entry.}

\medskip\noindent
{\bf Abstract.}

No abstract is available on the
official site as of November 20, 2026.

\section{Title to be announced: Sara Torelli}
\label{section-tba-torelli}
%% speaker: Sara Torelli, Politecnico di
%% speaker: Milano
%% place: LEGAL 2027, Sítio Assunção,
%% place: Teresópolis, RJ, Brazil

\medskip\noindent
{\bf Official presentation type.}

Invited speaker.

\medskip\noindent
{\bf Title.}

The official site lists the title as
``TBA'' as of November 20, 2026.

\medskip\noindent
{\bf Speaker.}

Sara Torelli, Politecnico di Milano.

\medskip\noindent
{\bf Schedule.}

Not listed on the official site as of
November 20, 2026.

\medskip\noindent
{\bf Official listing.}

\href{%
https://sites.google.com/view/%
geoalgcompluff/legal/%
legal-2027-resumos#h.59gubrdmgdxn%
}{Programme entry.}

\medskip\noindent
{\bf Abstract.}

No abstract is available on the
official site as of November 20, 2026.

\section{Title to be announced: Simone
Marchesi}
\label{section-tba-marchesi}
%% speaker: Simone Marchesi, Universidad de
%% speaker: Barcelona
%% place: LEGAL 2027, Sítio Assunção,
%% place: Teresópolis, RJ, Brazil

\medskip\noindent
{\bf Official presentation type.}

Invited speaker.

\medskip\noindent
{\bf Title.}

The official site lists the title as
``TBA'' as of November 20, 2026.

\medskip\noindent
{\bf Speaker.}

Simone Marchesi, Universidad de Barcelona.

\medskip\noindent
{\bf Schedule.}

Not listed on the official site as of
November 20, 2026.

\medskip\noindent
{\bf Official listing.}

\href{%
https://sites.google.com/view/%
geoalgcompluff/legal/%
legal-2027-resumos#h.59gubrdmgdxn%
}{Programme entry.}

\medskip\noindent
{\bf Abstract.}

No abstract is available on the
official site as of November 20, 2026.

\section{Rank 3 curve semistable Higgs
bundles}
\label{section-rank-3-curve-semistable-higgs}
%% speaker: Ugo Bruzzo, UFMG
%% place: LEGAL 2027, Sítio Assunção,
%% place: Teresópolis, RJ, Brazil

\medskip\noindent
{\bf Official presentation type.}

Invited speaker.

\medskip\noindent
{\bf Speaker.}

Ugo Bruzzo, UFMG.

\medskip\noindent
{\bf Schedule.}

Not listed on the official site as of
November 20, 2026.

\medskip\noindent
{\bf Official listing.}

\href{%
https://sites.google.com/view/%
geoalgcompluff/legal/%
legal-2027-resumos#h.59gubrdmgdxn%
}{Programme entry.}

\medskip\noindent
{\bf Abstract.}

Higgs bundles are vector bundles on a complex
projective variety equipped with an extra
twisted endomorphism, the Higgs field; they
arise in nonabelian Hodge theory and encode
representations of the fundamental group. A
natural way to test the stability of such an
object is to restrict it to all the curves
lying in the variety: a Higgs bundle is
called curve semistable if every such
restriction is semistable. A conjecture from
2011 predicts that this property is
equivalent to a numerical condition, the
vanishing of a characteristic class called
the discriminant, together with ordinary
semistability. One half of the statement is
known, and the conjecture holds for bundles
of rank two, but Huang has recently found a
counterexample in rank four. We show that the
conjecture is true for Higgs bundles of rank
three whose Higgs field is nilpotent, on
varieties of any dimension; as a consequence
it holds for all rank three Higgs bundles on
a large class of varieties, including smooth
surfaces in projective three-space. The proof
separates two cases according to how
degenerate the Higgs field is, and uses the
geometry of the Higgs Grassmannians, which
parametrize the subbundles preserved by the
Higgs field, together with a classical result
of Bogomolov on one-forms on surfaces.
Motivated by the counterexample, we then
propose a refined version of the conjecture:
instead of a single Higgs bundle, one should
look at the whole family obtained by
rescaling the Higgs field and at its limit,
the Hodge limit, inside Simpson’s moduli
space. We prove that this refined conjecture
holds in rank two and, thanks to the
nilpotent result, in rank three.

\section{Ehrhart theory for toric vector
bundles}
\label{section-ehrhart-theory-for-toric}
%% speaker: Ana Maria Botero, University of
%% speaker: Bielefeld
%% place: LEGAL 2027, Sítio Assunção,
%% place: Teresópolis, RJ, Brazil

\medskip\noindent
{\bf Official presentation type.}

Special session (mentor).

\medskip\noindent
{\bf Speaker.}

Ana Maria Botero, University of Bielefeld.

\medskip\noindent
{\bf Schedule.}

Not listed on the official site as of
November 20, 2026.

\medskip\noindent
{\bf Official listing.}

\href{%
https://sites.google.com/view/%
geoalgcompluff/legal/%
legal-2027-resumos#h.6f5k7zdb7c47%
}{Programme entry.}

\medskip\noindent
{\bf Abstract.}

One of the remarkable features of toric
geometry is that many geometric questions can
be translated into combinatorial problems
involving lattice polytopes. For toric line
bundles, this correspondence is particularly
elegant: global sections are described by
lattice points of a polytope, and their
Hilbert functions are governed by Ehrhart
theory. What happens when we replace line
bundles by vector bundles of higher rank?

Toric vector bundles admit a rich
combinatorial description in terms of
filtrations, matroids, and collections of
polytopes known as parliaments of polytopes.
These structures provide a way to study
spaces of global sections and suggest natural
generalizations of classical Ehrhart theory.

In this mentoring project, we will
investigate how the geometry of toric vector
bundles is reflected in the combinatorics of
their associated polytopes. Our guiding
questions will concern the behavior of
Hilbert functions and the positivity
properties of their coefficients: To what
extent do classical results from Ehrhart
theory extend to higher-rank vector bundles?
Can geometric positivity conditions on vector
bundles be characterized or detected through
their associated polytopes?

\section{Esquemas de Mustafin e
Grassmannianas de Aljavas}
\label{section-esquemas-de-mustafin-e}
%% speaker: Eduardo Esteves, IMPA
%% place: LEGAL 2027, Sítio Assunção,
%% place: Teresópolis, RJ, Brazil

\medskip\noindent
{\bf Official presentation type.}

Special session (mentor).

\medskip\noindent
{\bf Speaker.}

Eduardo Esteves, IMPA.

\medskip\noindent
{\bf Schedule.}

Not listed on the official site as of
November 20, 2026.

\medskip\noindent
{\bf Official listing.}

\href{%
https://sites.google.com/view/%
geoalgcompluff/legal/%
legal-2027-resumos#h.6f5k7zdb7c47%
}{Programme entry.}

\medskip\noindent
{\bf Abstract.}

O objetivo é uma introdução a Esquemas de
Mustafin e Grassmannianas de Aljavas, fazendo
um passeio pela Geometria Algébrica,
Geometria Tropical, Combinatória e Teoria de
Representações.

Esquemas de Mustafin são certas degenerações
de espaços projetivos. Mais especificamente,
surgem ao se considerar a diagonal pequena em
um produto de espaços projetivos de mesma
dimensão, e sua órbita pela ação dos grupos
de automorfismos dos fatores no espaço de
Hilbert correspondente. Variedades de
Mustafin são subesquemas do produto
parametrizados por um ponto no fecho da
órbita.

Esquemas de Mustafin são reduzidos e
Cohen–Macaulay. Possuem em geral várias
componentes irredutíveis. Sob certas
condições, pode-se mostrar que um esquema de
Mustafin é ”simple normal crossings,” o que
significa que suas componentes irredutíveis
são lisas, e as interseções destas
componentes são transversais (e lisas). Mais
ainda, neste caso, o grafo de interseções é o
grafo de adjacências dos reticulados que
descrevem a degeneração, e o esquema de
Mustafin é uma Grassmanniana de aljavas de
subespaços de dimensão 1.

Por outro lado, certas Grassmannianas de
aljavas vem sendo estudadas no contexto de
degenerações de séries lineares em famílias
de variedades algébricas. São chamadas de
”linked projective spaces”, espaços
projetivos ligados. Quando elas surgem de uma
tal degeneração, parametrizam os limites de
divisores associados à degeneração no esquema
de Hilbert da variedade limite. Mais ainda,
são esquemas de Mustafin.

Mas espaços projetivos ligados podem ser
considerados de forma independente. Uma das
principais questões na área é se eles são
sempre degenerações de espaços projetivos,
mais precisamente, se são esquemas de
Mustafin. Mostrou-se recentemente que têm a
classe de Chow e o polinômio de Hilbert
esperados de esquemas de Mustafin, os da
diagonal, e que são ”simple normal
crossings”.

A principal questão continua em aberto, mas,
curiosamente, mostrou-se a recíproca: toda
variedade de Mustafin é um espaço projetivo
ligado. Além dessa questão, muitas outras
surgem, e a ideia é entendê-las. As
principais referências para uma introdução à
teoria são [ESV26b, ESV26a, CHSW11]. Mais
referências, inclusive sobre resultados
mencionados acima, são [CS10, Li18, HL21,
HZ23, HZ25, Fal01, dLE26, dLEV26, dLEUar,
ERV25]

References

[CHSW11] Dustin Cartwright, Mathias H\"abich,
Bernd Sturmfels, and Annette Werner, Mustafin
varieties, Selecta Math. (N.S.) 17 (2011),
no. 4, 757–793.

[CS10] Dustin Cartwright and Bernd Sturmfels,
The Hilbert scheme of the diagonal in a
product of projective spaces, Int. Math. Res.
Not. IMRN 2010 (2010), no. 9, 1741–1771.

[dLE26] Felipe de Leon and Eduardo Esteves,
Polymatroidal tilings and the Chow class of
linked projective spaces, Bull. Lond. Math.
Soc. 58 (2026), no. 5, Paper No. e70370, 23.

[dLEUar] Felipe de León, Eduardo Esteves, and
Carlos Urtecho, Mustafin varieties and linked
nets, To appear.

[dLEV26] Felipe de León, Eduardo Esteves, and
Eduardo Vital, On the Hilbert polynomial of
the linked projective space, 2026.

[ERV25] Eduardo Esteves, Piere Rodr\'\i guez,
and Eduardo Vital, Degenerations of maps to
projective spaces, 2025.

[ESV26a] Eduardo Esteves, Renan Santos, and
Eduardo Vital, Quiver representations arising
from degenerations of linear series, II,
Bull. Braz. Math. Soc. (N.S.) 57 (2026), no.
1, Paper No. 4, 40.

[ESV26b] , Quiver representations arising
from degenerations of linear series, I, Comm.
Algebra 54 (2026), no. 2, 525–554.

[Fal01] Gerd Faltings, Toroidal resolutions
for some matrix singularities, Moduli of
abelian varieties (Texel Island, 1999),
Progr. Math., vol. 195, Birkh\"auser, Basel,
2001, pp. 157–184.

[HL21] Marvin Anas Hahn and Binglin Li,
Mustafin varieties, moduli spaces and
tropical geometry, Manuscripta Math. 166
(2021), no. 1-2, 167–197.

[HZ23] Xiang He and Naizhen Zhang,
Degenerations of Grassmannians via lattice
configurations, Int. Math. Res. Not. IMRN
(2023), no. 1, 298–349.

[HZ25] , Degenerations of Grassmannians via
lattice configurations II, Math. Ann. 391
(2025), no. 1, 965–1008.

[Li18] Binglin Li, Images of rational maps of
projective spaces, Int. Math. Res. Not. IMRN
2018 (2018), no. 13, 4190–4228.

\section{Title to be announced: Ethan
Cotterill}
\label{section-tba-cotterill}
%% speaker: Ethan Cotterill, UNICAMP
%% place: LEGAL 2027, Sítio Assunção,
%% place: Teresópolis, RJ, Brazil

\medskip\noindent
{\bf Official presentation type.}

Special session (mentor).

\medskip\noindent
{\bf Title.}

The official site lists the title as
``TBA'' as of November 20, 2026.

\medskip\noindent
{\bf Speaker.}

Ethan Cotterill, UNICAMP.

\medskip\noindent
{\bf Schedule.}

Not listed on the official site as of
November 20, 2026.

\medskip\noindent
{\bf Official listing.}

\href{%
https://sites.google.com/view/%
geoalgcompluff/legal/%
legal-2027-resumos#h.6f5k7zdb7c47%
}{Programme entry.}

\medskip\noindent
{\bf Abstract.}

No abstract is available on the
official site as of November 20, 2026.

\section{Um século de pontos racionais em
curvas algébricas: da tese de Artin aos dIAs
de hoje}
\label{section-um-seculo-de-pontos-racionais}
%% speaker: Herivelto Borges, USP
%% place: LEGAL 2027, Sítio Assunção,
%% place: Teresópolis, RJ, Brazil

\medskip\noindent
{\bf Official presentation type.}

Special session (mentor).

\medskip\noindent
{\bf Speaker.}

Herivelto Borges, USP.

\medskip\noindent
{\bf Schedule.}

Not listed on the official site as of
November 20, 2026.

\medskip\noindent
{\bf Official listing.}

\href{%
https://sites.google.com/view/%
geoalgcompluff/legal/%
legal-2027-resumos#h.6f5k7zdb7c47%
}{Programme entry.}

\medskip\noindent
{\bf Abstract.}

O curso trata do problema de estimar o número
de pontos racionais de uma curva de gênero
$g$ sobre $\mathbb F_q$. Começamos pelo
desenvolvimento histórico da hipótese de
Riemann para curvas, de Artin e Hasse até a
demonstração de Weil, e pela função zeta,
cujas propriedades acessíveis por métodos
elementares são demonstradas e da qual se
deduz a cota de Hasse–Weil. Em seguida,
estudamos os refinamentos posteriores: a cota
de Serre, obtida da integralidade dos
autovalores de Frobenius; as cotas de Ihara e
de Oesterlé, obtidas das contagens de pontos
sobre extensões; e a cota de Stöhr–Voloch, de
natureza geométrica, obtida da ação do mapa
de Frobenius e de suas relações com os
espaços osculadores. Discutimos o alcance de
cada uma dessas cotas e suas generalizações.A
segunda parte é dedicada ao estado atual da
teoria: curvas Frobenius não clássicas,
curvas maximais e a curva Hermitiana, como
exemplos em que as cotas são atingidas;
avanços recentes; e problemas em aberto, como
a determinação de $N_q(g)$, a classificação
de curvas maximais e questões assintóticas.
Se o tempo permitir, abordaremos conexões com
códigos corretores de erros, torres de corpos
de funções e geometria finita.

\section{Title to be announced: Karl Christ}
\label{section-tba-christ}
%% speaker: Karl Christ, Università degli
%% speaker: Studi di Torino
%% place: LEGAL 2027, Sítio Assunção,
%% place: Teresópolis, RJ, Brazil

\medskip\noindent
{\bf Official presentation type.}

Special session (mentor).

\medskip\noindent
{\bf Title.}

The official site lists the title as
``TBA'' as of November 20, 2026.

\medskip\noindent
{\bf Speaker.}

Karl Christ, Università degli Studi di
Torino.

\medskip\noindent
{\bf Schedule.}

Not listed on the official site as of
November 20, 2026.

\medskip\noindent
{\bf Official listing.}

\href{%
https://sites.google.com/view/%
geoalgcompluff/legal/%
legal-2027-resumos#h.6f5k7zdb7c47%
}{Programme entry.}

\medskip\noindent
{\bf Abstract.}

No abstract is available on the
official site as of November 20, 2026.

\section{Motivc classes of moduli spaces}
\label{section-motivc-classes-of-moduli}
%% speaker: Michele Graffeo, SISSA
%% place: LEGAL 2027, Sítio Assunção,
%% place: Teresópolis, RJ, Brazil

\medskip\noindent
{\bf Official presentation type.}

Special session (mentor).

\medskip\noindent
{\bf Speaker.}

Michele Graffeo, SISSA.

The main LEGAL page of the official site
lists the affiliation as Politecnico di
Milano.

\medskip\noindent
{\bf Schedule.}

Not listed on the official site as of
November 20, 2026.

\medskip\noindent
{\bf Official listing.}

\href{%
https://sites.google.com/view/%
geoalgcompluff/legal/%
legal-2027-resumos#h.6f5k7zdb7c47%
}{Programme entry.}

\medskip\noindent
{\bf Abstract.}

The study of moduli spaces and their geometry
is one of the central themes in algebraic
geometry. However, owing to their abstract
definition, geometric information about these
spaces is often inaccessible to direct
investigation. This motivates the
introduction of numerical invariants
associated with moduli spaces, which capture
and reflect some of their geometric
properties.

One of the most classical examples is the
Euler characteristic $\chi(X)\in\mathbb{Z}$,
an invariant associated with a sufficiently
well-behaved topological space $X$, which can
be used to distinguish spaces with different
topological properties. When dealing with
moduli spaces that are schemes of finite type
over $\mathbb{C}$, the required
"well-behavior" condition is satisfied,
making it meaningful to study such invariants
in order to gain insight into the geometry of
the underlying varieties.

A particularly remarkable feature of the
Euler characteristic is its compatibility
with two fundamental operations on varieties:
excision and Cartesian products. More
precisely, if $Y\subset X$ is a closed
subvariety, then $\chi(X)=\chi(X\setminus
Y)+\chi(Y)$, while for two varieties $X$ and
$Y$, we have $\chi(X\times
Y)=\chi(X)\cdot\chi(Y)$.

The Euler characteristic thus provides a
first example of a motivic invariant: a
homomorphism from the free abelian group
generated by isomorphism classes of complex
algebraic varieties to a ring, compatible
with excision and products. This naturally
leads to the notion of the Grothendieck ring
of varieties
$K_0(\operatorname{Var}_{\mathbb{C}})$. It is
the quotient of the polynomial ring over
$\mathbb{Z}$ generated by the isomorphism
classes $[X]$ of complex algebraic varieties
$X$ by the ideal generated by the elements
$$
[X]-[X\setminus Y]-[Y]
$$
and
$$
[X\times Y]-[X]\cdot[Y],
$$
where $Y\subset X$ ranges over closed
subvarieties. The elements of
$K_0(\operatorname{Var}_{\mathbb{C}})$ are
called motivic classes or simply motives. In
this framework, a motivic invariant is a ring
homomorphism
$K_0(\operatorname{Var}_{\mathbb{C}})
\longrightarrow A$, where $A$ is a
commutative ring. In particular, knowing the
motivic class of a variety provides access to
all motivic invariants associated with it.
Examples include the Euler characteristic and
the Hodge--Deligne polynomial.

Despite its elementary definition, the
Grothendieck ring of varieties exhibits a
number of striking pathologies. For instance,
the class of the affine line
$\mathbb{L}:=[\mathbb{A}^1]\in
K_0(\operatorname{Var}_{\mathbb{C}}),$ known
as the Lefschetz motive, is a zero-divisor,
[1]. In this setting, expressing the class
$[X]\in K_0(\operatorname{Var}_{\mathbb{C}})$
of a variety $X$ in terms of the classes of
simpler varieties can therefore provide
valuable insight into its geometry and reveal
relations that are not apparent from its
definition alone.

In enumerative geometry, it is often fruitful
to study families of moduli spaces indexed by
a discrete invariant, with the aim of
understanding their behaviour as the
invariant varies. A fundamental example is
provided by the Hilbert scheme
$\operatorname{Hilb}^d\mathbb{A}^n$ of $d$
points in affine $n$-space, which
parametrises zero-dimensional closed
subschemes of $\mathbb{A}^n$ of length $d$,
for some $n,d\ge 0$. In this context, one is
naturally led to consider generating series
of the form

$$
\mathsf{H}_n(t) =\sum_{d\geq
0}[\operatorname{Hilb}^d\mathbb{A}^n]t^d,
\qquad \mathsf{Z}_d(t) =\sum_{n\geq
0}[\operatorname{Hilb}^d\mathbb{A}^n]t^n,
$$

both taking values in
$K_0(\operatorname{Var}_{\mathbb{C}})[\![
t]\!]$. Remarkably, while the motivic
generating series $\mathsf{H}_n(t)$ remains
largely mysterious for $n\ge 3$, the series
$\mathsf{Z}_d(t)$ has recently been
investigated, leading to explicit expressions
in terms of rational functions and finitely
many initial data [6]. These results
highlight the rich interplay between the
geometry of Hilbert schemes, combinatorics
and the theory of motivic invariants.
Moreover, knowledge of the behaviour of the
series $\mathsf{Z}_d(t)$, for $d\geq 0$, has
also made it possible to describe certain
properties of the ind-varieties
$\operatorname{Hilb}^d\mathbb{A}^{\infty}
:=\varinjlim_{n}\operatorname{Hilb}^d
\mathbb{A} ^n.$ It is also worth mentioning
that knowledge of the motivic classes
$[\operatorname{Hilb}^d\mathbb{A}^n]$ in some
range $0\le d\le d'$ and $0\le n\le n'$ makes
it possible to compute the motivic classes of
$\operatorname{Hilb}^d X$ for smooth
varieties $X$ of dimensions $n$ in the same
range.

It has been shown in several works that the
presence of a torus action provides access to
the powerful tool of Bia{\l}ynicki-Birula
decomposition, which can be used to compute
motivic classes [2, 3, 5]. Recent results
obtained in [6, 4] rely primarily on the
interplay between Bia{\l}ynicki-Birula
decomposition and algebraic techniques such
as \textit{Macaulay duality} also known as
\textit{apolarity}. Indeed, combining these
approaches makes it possible to break down
the computation of motivic classes of
varieties equipped with a torus action into
more manageable subproblems.

As mentioned above, this approach applies
more generally to varieties equipped with an
algebraic torus action. Classical examples
include Quot schemes and nested versions of
Hilbert and Quot schemes. For many of these
families of moduli spaces, the structure of
the generating series of their motivic
classes remains unknown, leaving ample room
for further investigation in this direction.

References:

1. Lev~A. Borisov. - The class of the affine
line is a zero divisor in the {G}rothendieck
ring - J. Algebraic Geom., 27(2):203--209,
2018.

2. Jan Cheah. - The cohomology of smooth
nested {H}ilbert schemes of points. -
ProQuest LLC, Ann Arbor, MI, 1994. - Thesis
(Ph.D.)--The University of Chicago.

3. Jan Cheah. - On the cohomology of
{H}ilbert schemes of points. - J. Algebraic
Geom., 5(3):479--511, 1996.

4. Nadir Fasola, Michele Graffeo, Danilo
Lewański, and Andrea~T. Ricolfi. - Invariants
of nested Hilbert and Quot schemes on
surfaces. - J. Math. Pures Appl. (9),
216:103973, 2026.

5. Lothar G{\"{o}}ttsche. - On the motive of
the Hilbert scheme of points on a surface. -
Math. Res. Lett., 8:613--627, 2001.

6. Michele Graffeo, Sergej Monavari, Riccardo
Moschetti, and Andrea~T. Ricolfi. - The
motive of the Hilbert scheme of points in all
dimensions. - Proc. Lond. Math. Soc. (3),
132(3):e70140, 2026.

\section{Title to be announced: Olivier
Martin}
\label{section-tba-martin}
%% speaker: Olivier Martin, IMPA
%% place: LEGAL 2027, Sítio Assunção,
%% place: Teresópolis, RJ, Brazil

\medskip\noindent
{\bf Official presentation type.}

Special session (mentor).

\medskip\noindent
{\bf Title.}

The official site lists the title as
``TBA'' as of November 20, 2026.

\medskip\noindent
{\bf Speaker.}

Olivier Martin, IMPA.

\medskip\noindent
{\bf Schedule.}

Not listed on the official site as of
November 20, 2026.

\medskip\noindent
{\bf Official listing.}

\href{%
https://sites.google.com/view/%
geoalgcompluff/legal/%
legal-2027-resumos#h.6f5k7zdb7c47%
}{Programme entry.}

\medskip\noindent
{\bf Abstract.}

No abstract is available on the
official site as of November 20, 2026.

\section{Compact Hyperkähler Manifolds}
\label{section-compact-hyperkahler-manifolds}
%% speaker: Yulieth Prieto, Universidad
%% speaker: Católica de Chile
%% place: LEGAL 2027, Sítio Assunção,
%% place: Teresópolis, RJ, Brazil

\medskip\noindent
{\bf Official presentation type.}

Special session (mentor).

\medskip\noindent
{\bf Speaker.}

Yulieth Prieto, Universidad Católica de
Chile.

\medskip\noindent
{\bf Schedule.}

Not listed on the official site as of
November 20, 2026.

\medskip\noindent
{\bf Official listing.}

\href{%
https://sites.google.com/view/%
geoalgcompluff/legal/%
legal-2027-resumos#h.6f5k7zdb7c47%
}{Programme entry.}

\medskip\noindent
{\bf Abstract.}

Compact hyperkähler (HK) manifolds, i.e.
simply connected compact Kähler manifolds X
with $H^0(X,\Omega^2_X)$ spanned by an
everywhere non-degenerate holomorphic 2-form,
are one of the three building blocks of
Kähler manifolds with trivial first Chern
class (Beauville--Bogomolov decomposition).
They are higher-dimensional analogues of K3
surfaces and enjoy a remarkably rich
geometry, controlled to a large extent by the
Hodge structure on $H^2(X,Z)$ endowed with
the Beauville--Bogomolov--Fujiki form. The
course is an invitation to the subject
through examples: we will see how HK
manifolds arise from K3 surfaces and from
some Fano varieties; how they shed light on
classical questions such as the rationality
of cubic fourfolds; and why producing new
examples is a genuinely hard problem.

\section{Folheações e matrizes 2 x (n +1)}
\label{section-folheacoes-e-matrizes-2-x-n-1}
%% speaker: Alan Muniz, UFPE
%% place: LEGAL 2027, Sítio Assunção,
%% place: Teresópolis, RJ, Brazil

\medskip\noindent
{\bf Official presentation type.}

Contributed talk.

\medskip\noindent
{\bf Speaker.}

Alan Muniz, UFPE.

\medskip\noindent
{\bf Schedule.}

Not listed on the official site as of
November 20, 2026.

\medskip\noindent
{\bf Official listing.}

\href{%
https://sites.google.com/view/%
geoalgcompluff/legal/%
legal-2027-resumos#h.dant8g3oe5ea%
}{Programme entry.}

\medskip\noindent
{\bf Abstract.}

Dado um par de polinômios homogêneos coprimos
(f, g) temos uma folheação racional cujo
feixe tangente pode ser determinado como
núcleo da matriz jacobiana de (f, g). Para
uma matriz A, 2 x (n+1), com linhas
homogêneas, não necessariamente jacobiana,
seu núcleo determina uma distribuição, em
geral, não integrável. Uma pergunta natural é
se a integrabilidade implica que a folheação
seja racional. Vamos discutir quando isso
vale e como se relaciona com uma conjectura
de Mendson e Pereira.

\section{Uma visita ao Problema de Waring e
as Propriedades de Lefschetz}
\label{section-uma-visita-ao-problema-de}
%% speaker: Alejandro Vanegas Obregón,
%% speaker: Universidade Federal de Minas
%% speaker: Gerais
%% place: LEGAL 2027, Sítio Assunção,
%% place: Teresópolis, RJ, Brazil

\medskip\noindent
{\bf Official presentation type.}

Contributed talk.

\medskip\noindent
{\bf Speaker.}

Alejandro Vanegas Obregón, Universidade
Federal de Minas Gerais.

\medskip\noindent
{\bf Schedule.}

Not listed on the official site as of
November 20, 2026.

\medskip\noindent
{\bf Official listing.}

\href{%
https://sites.google.com/view/%
geoalgcompluff/legal/%
legal-2027-resumos#h.dant8g3oe5ea%
}{Programme entry.}

\medskip\noindent
{\bf Abstract.}

O trabalho desenvolve três variações do
problema de Waring para polinômios, focando
no posto de Waring, posto borda e posto de
cactus de uma forma homogênea. A ideia
central é analisar a influência das
propriedades de Lefschetz (tanto WLP quanto
SLP) da álgebra graduada associada nos postos
mencionados. Usando a teoria de matrizes
Hessianas mistas, tem-se conexões entre a
geometria das variedades secantes e as
propriedades das álgebras Artinianas de
Gorenstein. Isto deriva na construção de
novas famílias de formas silvestres (ou wild
forms) — definidas como formas cujo posto de
cactus, de natureza esquemática, é
estritamente maior do que o seu posto borda,
definido geométricamente. Se analisam também
formas com Hessiano nulo e se estabelecem
novos critérios para determinar postos em
contextos onde os métodos clássicos são
insuficientes. Keywords: Problema de Waring,
Propriedades de Lefschetz, Matriz Hessiana,
Formas Silvestres.

\section{Exceptional collections on canonical
stacks of log del Pezzo surfaces}
\label{section-exceptional-collections-on}
%% speaker: Alex Junior Gomez Saltachin,
%% speaker: PUC-Rio
%% place: LEGAL 2027, Sítio Assunção,
%% place: Teresópolis, RJ, Brazil

\medskip\noindent
{\bf Official presentation type.}

Contributed talk.

\medskip\noindent
{\bf Speaker.}

Alex Junior Gomez Saltachin, PUC-Rio.

\medskip\noindent
{\bf Schedule.}

Not listed on the official site as of
November 20, 2026.

\medskip\noindent
{\bf Official listing.}

\href{%
https://sites.google.com/view/%
geoalgcompluff/legal/%
legal-2027-resumos#h.dant8g3oe5ea%
}{Programme entry.}

\medskip\noindent
{\bf Abstract.}

We study the derived categories of canonical
smooth Deligne–Mumford stacks associated with
log del Pezzo surfaces with singularities of
type 1/3 (1,1). Using the geometry of log del
Pezzo cascades together with the special
McKay correspondence and Orlov’s blow-up
formula, we give an explicit construction of
full exceptional collections for these
canonical stacks. As a concrete example, we
consider the sporadic Johnson–Kollár surface
$X_{10}\subset P(1,2,3,5)$.

We describe its geometry in terms of the
blow-up of $P(1,1,3)$ at eight general points
and obtain a full exceptional collection of
length 13 on its canonical stack. If time
permits, I will also discuss how this picture
is related to semiorthogonal decompositions
of the singular surface itself.

\section{On automorphisms of hypersurfaces in
Fano and Calabi Yau varieties}
\label{section-on-automorphisms-of}
%% speaker: Ana Victoria Martins Quedo,
%% speaker: Università degli studi di Ferrara
%% place: LEGAL 2027, Sítio Assunção,
%% place: Teresópolis, RJ, Brazil

\medskip\noindent
{\bf Official presentation type.}

Contributed talk.

\medskip\noindent
{\bf Speaker.}

Ana Victoria Martins Quedo, Università degli
studi di Ferrara.

\medskip\noindent
{\bf Schedule.}

Not listed on the official site as of
November 20, 2026.

\medskip\noindent
{\bf Official listing.}

\href{%
https://sites.google.com/view/%
geoalgcompluff/legal/%
legal-2027-resumos#h.dant8g3oe5ea%
}{Programme entry.}

\medskip\noindent
{\bf Abstract.}

The classical theorem of Matsumura–Monsky
states that every automorphism of a smooth
hypersurface $X \subseteq \mathbb{P}^n$ of
degree d is induced by an automorphism of the
ambient projective space for $d \geq 2$ and
$n \geq 3$, with the sole exception of smooth
quartic surfaces in $\mathbb{P}^3$ . More
recently, Esser established an analogous
result for quasismooth hypersurfaces in
well-formed weighted projective spaces. One
can ask whether analogous results still hold
for more general ambient spaces, such as Fano
or Calabi–Yau manifolds. In this talk, we
present work in progress that provides a
partial answer to this question, relating the
above problem to more modern concepts in
algebraic geometry and commutative algebra
such as normally presented line bundles.

\section{On the Theta Divisor and
Torelli-like Theorems in the Tropical
Context}
\label{section-on-the-theta-divisor-and}
%% speaker: Analice Antonella Silva Guedes,
%% speaker: Universidade Federal Fluminense
%% place: LEGAL 2027, Sítio Assunção,
%% place: Teresópolis, RJ, Brazil

\medskip\noindent
{\bf Official presentation type.}

Contributed talk.

\medskip\noindent
{\bf Speaker.}

Analice Antonella Silva Guedes, Universidade
Federal Fluminense.

\medskip\noindent
{\bf Schedule.}

Not listed on the official site as of
November 20, 2026.

\medskip\noindent
{\bf Official listing.}

\href{%
https://sites.google.com/view/%
geoalgcompluff/legal/%
legal-2027-resumos#h.dant8g3oe5ea%
}{Programme entry.}

\medskip\noindent
{\bf Abstract.}

We investigate properties of posets of
quasistable pseudodivisors on a given
pointed, undirected, connected multigraph as
an analogue of theta divisor for pointed
algebraic curves. In particular, to which
extent one can recover the base graph from
the datum of the corresponding poset.

\section{Nikulin configurations on Kummer
surfaces}
\label{section-nikulin-configurations-on}
%% speaker: Carlos Daniel Rivera Chacón,
%% speaker: University of La Frontera
%% place: LEGAL 2027, Sítio Assunção,
%% place: Teresópolis, RJ, Brazil

\medskip\noindent
{\bf Official presentation type.}

Contributed talk.

\medskip\noindent
{\bf Speaker.}

Carlos Daniel Rivera Chacón, University of La
Frontera.

\medskip\noindent
{\bf Schedule.}

Not listed on the official site as of
November 20, 2026.

\medskip\noindent
{\bf Official listing.}

\href{%
https://sites.google.com/view/%
geoalgcompluff/legal/%
legal-2027-resumos#h.dant8g3oe5ea%
}{Programme entry.}

\medskip\noindent
{\bf Abstract.}

The official site lists the abstract
as ``TBA'' as of November 20, 2026.

\section{Smoothing of a Stanley-Reisner
variety}
\label{section-smoothing-of-a-stanley}
%% speaker: Daniel González Casanova Azuela,
%% speaker: PUC-Rio
%% place: LEGAL 2027, Sítio Assunção,
%% place: Teresópolis, RJ, Brazil

\medskip\noindent
{\bf Official presentation type.}

Contributed talk.

\medskip\noindent
{\bf Speaker.}

Daniel González Casanova Azuela, PUC-Rio.

\medskip\noindent
{\bf Schedule.}

Not listed on the official site as of
November 20, 2026.

\medskip\noindent
{\bf Official listing.}

\href{%
https://sites.google.com/view/%
geoalgcompluff/legal/%
legal-2027-resumos#h.dant8g3oe5ea%
}{Programme entry.}

\medskip\noindent
{\bf Abstract.}

In a 1967 paper, Grünbaum and Sreedharan
discovered a triangulation M of the 3-sphere
with interesting combinatorial properties.
The Stanley-Reisner variety SR(M) is a union
of linear subspaces in $P^7$ , and is highly
singular. We look for a smoothing of SR(M),
that is, a flat family whose special fibre is
SR(M) and whose general fibre is a smooth
variety.

I will present a preliminary AI-assisted
construction of the desired smoothing. The
procedure is, roughly, finding finite-order
deformations which altogether yield an
abstract object called ""formal
deformation"", and then ""algebraizing"" such
object to produce the desired family of
algebraic varieties.

\section{The universal tropical theta
divisor}
\label{section-the-universal-tropical-theta}
%% speaker: Giusi Capobianco, University of
%% speaker: Tor Vergata
%% place: LEGAL 2027, Sítio Assunção,
%% place: Teresópolis, RJ, Brazil

\medskip\noindent
{\bf Official presentation type.}

Contributed talk.

\medskip\noindent
{\bf Speaker.}

Giusi Capobianco, University of Tor Vergata.

\medskip\noindent
{\bf Schedule.}

Not listed on the official site as of
November 20, 2026.

\medskip\noindent
{\bf Official listing.}

\href{%
https://sites.google.com/view/%
geoalgcompluff/legal/%
legal-2027-resumos#h.dant8g3oe5ea%
}{Programme entry.}

\medskip\noindent
{\bf Abstract.}

In joint work with A. Guedes and M. Pacini we
study the geometric properties of the
tropical theta divisor in order to fully
describe its stratification. By using the
poset quasistable divisors we are able to
describe the polyhedral complex structure of
Theta and define a universal tropical theta
divisor as generalized cone complex.

\section{Non-existence of separating
morphisms of low degree}
\label{section-non-existence-of-separating}
%% speaker: Gurvan Mével, Sorbonne Université
%% speaker: - IMPA
%% place: LEGAL 2027, Sítio Assunção,
%% place: Teresópolis, RJ, Brazil

\medskip\noindent
{\bf Official presentation type.}

Contributed talk.

\medskip\noindent
{\bf Speaker.}

Gurvan Mével, Sorbonne Université - IMPA.

\medskip\noindent
{\bf Schedule.}

Not listed on the official site as of
November 20, 2026.

\medskip\noindent
{\bf Official listing.}

\href{%
https://sites.google.com/view/%
geoalgcompluff/legal/%
legal-2027-resumos#h.dant8g3oe5ea%
}{Programme entry.}

\medskip\noindent
{\bf Abstract.}

Given a real algebraic curve C, a morphism $C
\to P^1$ is separating if the preimage of
real points of $P^1$ is exactly the real
points of C. The degree of a separating
morphism is necessarily greater than the
number of connected components of the real
part of C, but this lower bound is not always
achieved. In this talk we will present
obstruction for the existence of separating
morphism of low degree. This is joint work
with A. Demory and A. Toussaint, based on
ideas by M. Manzaroli.

\section{Lefschetz Properties: A Geometric
Viewpoint}
\label{section-lefschetz-properties-a}
%% speaker: Janaíne Geralda Mesquita Martins,
%% speaker: Universidade Federal de Minas
%% speaker: Gerais
%% place: LEGAL 2027, Sítio Assunção,
%% place: Teresópolis, RJ, Brazil

\medskip\noindent
{\bf Official presentation type.}

Contributed talk.

\medskip\noindent
{\bf Speaker.}

Janaíne Geralda Mesquita Martins,
Universidade Federal de Minas Gerais.

\medskip\noindent
{\bf Schedule.}

Not listed on the official site as of
November 20, 2026.

\medskip\noindent
{\bf Official listing.}

\href{%
https://sites.google.com/view/%
geoalgcompluff/legal/%
legal-2027-resumos#h.dant8g3oe5ea%
}{Programme entry.}

\medskip\noindent
{\bf Abstract.}

The Lefschetz properties, inspired by the
classical Hard Lefschetz Theorem, have been
extensively studied in several areas of
mathematics, including algebra, topology,
combinatorics, and algebraic geometry. In
this talk, we will discuss the Lefschetz
properties of Artinian graded algebras from a
geometric viewpoint. In particular, we will
explore how geometric objects naturally
associated with these algebras can provide
information about the behavior of
multiplication by linear forms. We will
present some recent results illustrating
these connections, with particular emphasis
on geometric interpretations of the failure
of the Lefschetz properties. This work is
being carried out under the supervision of
and in collaboration with Prof. Rosa Maria
Miró-Roig (Universitat de Barcelona).

\section{Title to be announced: Joaquim Costa
Ribeiro}
\label{section-tba-ribeiro}
%% speaker: Joaquim Costa Ribeiro, IMPA
%% place: LEGAL 2027, Sítio Assunção,
%% place: Teresópolis, RJ, Brazil

\medskip\noindent
{\bf Official presentation type.}

Contributed talk.

\medskip\noindent
{\bf Title.}

The official site lists the title as
``TBA'' as of November 20, 2026.

\medskip\noindent
{\bf Speaker.}

Joaquim Costa Ribeiro, IMPA.

\medskip\noindent
{\bf Schedule.}

Not listed on the official site as of
November 20, 2026.

\medskip\noindent
{\bf Official listing.}

\href{%
https://sites.google.com/view/%
geoalgcompluff/legal/%
legal-2027-resumos#h.dant8g3oe5ea%
}{Programme entry.}

\medskip\noindent
{\bf Abstract.}

The official site lists the abstract
as ``TBA'' as of November 20, 2026.

\section{Equivalência racional de pontos em
hipersuperficies muito gerais de variedades
homogeneas}
\label{section-equivalencia-racional-de}
%% speaker: Juan Pablo Ocampo Gamboa, IMPA
%% place: LEGAL 2027, Sítio Assunção,
%% place: Teresópolis, RJ, Brazil

\medskip\noindent
{\bf Official presentation type.}

Contributed talk.

\medskip\noindent
{\bf Speaker.}

Juan Pablo Ocampo Gamboa, IMPA.

\medskip\noindent
{\bf Schedule.}

Not listed on the official site as of
November 20, 2026.

\medskip\noindent
{\bf Official listing.}

\href{%
https://sites.google.com/view/%
geoalgcompluff/legal/%
legal-2027-resumos#h.dant8g3oe5ea%
}{Programme entry.}

\medskip\noindent
{\bf Abstract.}

Seja X uma interseção completa muito geral de
multigrau $(d_1,\ldots,d_k)$ em uma variedade
homogênea $G/P$ de posto de Picard (1). No
caso $G/P=\mathbb{P}^{n+k}$, Voisin provou
que, se $\sum_i(d_i-1)>2n+1$, então quaisquer
dois pontos distintos de X não são
racionalmente equivalentes. Em 2021,
Chen–Lewis–Sheng melhoraram essa cota para
hipersuperfícies $k=1$, mostrando que
$(d_1-1>2n)$ é suficiente e ótimo.

Neste trabalho, buscamos estender esse
resultado para hipersuperfícies muito gerais
em variedades homogêneas $G/P$ de posto de
Picard (1), mostrando que, para $(d_1)$
suficientemente grande, quaisquer dois pontos
distintos de X não são racionalmente
equivalentes.

\section{Moduli Spaces of Rank 2 Torsion-Free
Sheaves with Quasi-Maximal Third Chern Class}
\label{section-moduli-spaces-of-rank-2}
%% speaker: Leonardo Silva de Oliveira, IMECC
%% speaker: - Unicamp
%% place: LEGAL 2027, Sítio Assunção,
%% place: Teresópolis, RJ, Brazil

\medskip\noindent
{\bf Official presentation type.}

Contributed talk.

\medskip\noindent
{\bf Speaker.}

Leonardo Silva de Oliveira, IMECC - Unicamp.

\medskip\noindent
{\bf Schedule.}

Not listed on the official site as of
November 20, 2026.

\medskip\noindent
{\bf Official listing.}

\href{%
https://sites.google.com/view/%
geoalgcompluff/legal/%
legal-2027-resumos#h.dant8g3oe5ea%
}{Programme entry.}

\medskip\noindent
{\bf Abstract.}

Hartshorne established the irreducibility and
smoothness of the Gieseker-Maruyama moduli
spaces parameterizing semistable rank 2
reflexive sheaves on $P^3$ with $c_1=-1$ and
the maximal third Chern class $c_3 = c_2^2$.
Later, Okonek and Spindler and Schmidt, on an
independent work, prove that this still true
for torsion-free sheaves. However, a gap
theorem by Miró-Roig shows that the moduli
space of reflexive sheaves is empty for a
certain range on the $c_3$, including the
quasi-maximal case $c_3 = c_2^2 - 2$. Despite
the absence of reflexive sheaves,
torsion-free sheaves do exist in this range,
Almeida, Jardim, and Tikhomirov constructed a
family of components for these spaces. In
this talk we will use the modular
Hartshorne-Serre correspondence to prove the
irreducibility of $M(-1, c_2, c_2^2 - 2)$ for
$c_2 >3$. (Joint work with Charles Almeida
and Marcos Jardim)

\section{Classification and Multiplicities of
Hecke Modifications on the Projective Line
over a Finite Field}
\label{section-classification-and}
%% speaker: Leonardo Soares Moço, Unicamp
%% place: LEGAL 2027, Sítio Assunção,
%% place: Teresópolis, RJ, Brazil

\medskip\noindent
{\bf Official presentation type.}

Contributed talk.

\medskip\noindent
{\bf Speaker.}

Leonardo Soares Moço, Unicamp.

\medskip\noindent
{\bf Schedule.}

Not listed on the official site as of
November 20, 2026.

\medskip\noindent
{\bf Official listing.}

\href{%
https://sites.google.com/view/%
geoalgcompluff/legal/%
legal-2027-resumos#h.dant8g3oe5ea%
}{Programme entry.}

\medskip\noindent
{\bf Abstract.}

We study Hecke modifications of vector
bundles on the projective line over a finite
field. Given a vector bundle E and a closed
point x, we ask which bundles E' can be
obtained through an elementary Hecke
modification of E at x, and with what
multiplicities each isomorphism class occurs.

The classification of vector bundles on $P^1$
reduces this problem to a combinatorial one
involving the degrees of the line bundle
summands. We use the Hall algebra of
Coh($P^1$) to describe the possible modified
bundles and to encode the corresponding
extension counts.

For degree-one points, we go further and
compute explicitly the multiplicity of each
possible modification. This provides a
concrete illustration of how Hall algebra
techniques can be used to study Hecke
modifications and extension problems for
vector bundles over finite fields.

\section{Vector bundles on stacky curves and
their moduli spaces}
\label{section-vector-bundles-on-stacky}
%% speaker: Lorenzo Cortelli, SISSA / IMPA
%% place: LEGAL 2027, Sítio Assunção,
%% place: Teresópolis, RJ, Brazil

\medskip\noindent
{\bf Official presentation type.}

Contributed talk.

\medskip\noindent
{\bf Speaker.}

Lorenzo Cortelli, SISSA / IMPA.

\medskip\noindent
{\bf Schedule.}

Not listed on the official site as of
November 20, 2026.

\medskip\noindent
{\bf Official listing.}

\href{%
https://sites.google.com/view/%
geoalgcompluff/legal/%
legal-2027-resumos#h.dant8g3oe5ea%
}{Programme entry.}

\medskip\noindent
{\bf Abstract.}

Vector bundles over a smooth projective curve
are among the central objects of algebraic
geometry. The corresponding moduli spaces
provide a remarkable class of algebraic
varieties, namely Fano varieties. After
recalling some of their fundamental
properties, including a classical theorem of
Narasimhan and Ramanan on their deformations,
we turn to the more general setting of stacky
curves.

Vector bundles on stacky curves naturally
incorporate representation-theoretic data at
marked points and generalise the theory of
parabolic bundles. We discuss the geometry of
the associated moduli spaces and present
recent results showing that, under suitable
hypotheses, their deformation theory remains
entirely governed by the underlying stacky
curve. In particular, we obtain a stacky
analogue of the Narasimhan–Ramanan theorem,
showing that every infinitesimal deformation
of the moduli space arises from a deformation
of the curve itself.

\section{A stacky nonabelian hodge
correspondence for open Curves}
\label{section-a-stacky-nonabelian-hodge}
%% speaker: Luca Morstabilini, SISSA
%% place: LEGAL 2027, Sítio Assunção,
%% place: Teresópolis, RJ, Brazil

\medskip\noindent
{\bf Official presentation type.}

Contributed talk.

\medskip\noindent
{\bf Speaker.}

Luca Morstabilini, SISSA.

\medskip\noindent
{\bf Schedule.}

Not listed on the official site as of
November 20, 2026.

\medskip\noindent
{\bf Official listing.}

\href{%
https://sites.google.com/view/%
geoalgcompluff/legal/%
legal-2027-resumos#h.dant8g3oe5ea%
}{Programme entry.}

\medskip\noindent
{\bf Abstract.}

Let X be a compact Riemann surface. In the
90s, Simpson showed that there is a
correspondence (the so-called Nonabelian
Hodge correspondence) between stable Higgs
bundles on X and representations of the
fundamental group of X, therefore comparing
objects which depend on the holomorphic
structure of X with others depending only on
the topology. Moreover, the nonabelian Hodge
isomorphism lifts to homeomorphisms of the
moduli spaces parametrizing the
aforementioned objects. It does not, however,
extend to an equivalence of the relevant
moduli stacks. A stacky version of the NAH
correspondence for compact curves has
recently been proposed and proven by Davison,
Hennecart, and Schlegel-Mejia using
techniques from representation theory. After
reviewing their work, I will explain how the
stacky NAH correspondence extends to open
curves. At the level of moduli spaces, this
was done by Simpson; at the stacky level,
this relies on the study of sheaves root
stacks.

\section{Pointed Hurwitz Schemes}
\label{section-pointed-hurwitz-schemes}
%% speaker: Manoel Oliveira, Universidade
%% speaker: Federal Fluminense - UFF
%% place: LEGAL 2027, Sítio Assunção,
%% place: Teresópolis, RJ, Brazil

\medskip\noindent
{\bf Official presentation type.}

Contributed talk.

\medskip\noindent
{\bf Speaker.}

Manoel Oliveira, Universidade Federal
Fluminense - UFF.

\medskip\noindent
{\bf Schedule.}

Not listed on the official site as of
November 20, 2026.

\medskip\noindent
{\bf Official listing.}

\href{%
https://sites.google.com/view/%
geoalgcompluff/legal/%
legal-2027-resumos#h.dant8g3oe5ea%
}{Programme entry.}

\medskip\noindent
{\bf Abstract.}

We introduce the notion of a pointed
admissible cover, which essentially consists
of a finite morphism $\pi:C \to B$ between
nodal algebraic curves, with C endowed with
marked points, such that the images of these
points under $\pi$, together with the branch
locus of $\pi$, endow B with the structure of
a stable pointed curve. As in the case of
stable pointed curves - where the underlying
curve need not itself be stable - the cover
$\pi$ is, in general, not admissible in the
classical sense of Harris-Mumford.

We show that the moduli problem for families
of pointed admissible covers admits a coarse
moduli space which is a scheme, and we
describe its relation to the Hurwitz scheme
$H_{k,b}$ by means of a forgetful morphism.
This morphism forgets the marked points of C
and induces on $\pi:C \to B$ a kind of
"admissibilization", in analogy with the
stabilization process for $\overline{M}_{g,n}
\to \overline{M}_{g}$.

\section{Explicit birational geometry of
determinantal quartic 3-folds}
\label{section-explicit-birational-geometry}
%% speaker: Manuel Leal, IMPA
%% place: LEGAL 2027, Sítio Assunção,
%% place: Teresópolis, RJ, Brazil

\medskip\noindent
{\bf Official presentation type.}

Contributed talk.

\medskip\noindent
{\bf Speaker.}

Manuel Leal, IMPA.

\medskip\noindent
{\bf Schedule.}

Not listed on the official site as of
November 20, 2026.

\medskip\noindent
{\bf Official listing.}

\href{%
https://sites.google.com/view/%
geoalgcompluff/legal/%
legal-2027-resumos#h.dant8g3oe5ea%
}{Programme entry.}

\medskip\noindent
{\bf Abstract.}

The official site lists the abstract
as ``TBA'' as of November 20, 2026.

\section{Classifying QHD Surfaces}
\label{section-classifying-qhd-surfaces}
%% speaker: Marcos Canedo, Pontificia
%% speaker: Universidad Católica de Chile
%% place: LEGAL 2027, Sítio Assunção,
%% place: Teresópolis, RJ, Brazil

\medskip\noindent
{\bf Official presentation type.}

Contributed talk.

\medskip\noindent
{\bf Speaker.}

Marcos Canedo, Pontificia Universidad
Católica de Chile.

\medskip\noindent
{\bf Schedule.}

Not listed on the official site as of
November 20, 2026.

\medskip\noindent
{\bf Official listing.}

\href{%
https://sites.google.com/view/%
geoalgcompluff/legal/%
legal-2027-resumos#h.dant8g3oe5ea%
}{Programme entry.}

\medskip\noindent
{\bf Abstract.}

I will present combinatorial techniques to
classify normal projective surfaces which are
limit surfaces of Q-Gorenstein degenerations
with constant Betti numbers, i.e. the Betti
numbers of the members of the family remains
constant under the degeneration.

\section{From Combinatorial to Algebraic
invariants via Compactifications}
\label{section-from-combinatorial-to}
%% speaker: Myrla Barbosa, UNEAL
%% place: LEGAL 2027, Sítio Assunção,
%% place: Teresópolis, RJ, Brazil

\medskip\noindent
{\bf Official presentation type.}

Contributed talk.

\medskip\noindent
{\bf Speaker.}

Myrla Barbosa, UNEAL.

\medskip\noindent
{\bf Schedule.}

Not listed on the official site as of
November 20, 2026.

\medskip\noindent
{\bf Official listing.}

\href{%
https://sites.google.com/view/%
geoalgcompluff/legal/%
legal-2027-resumos#h.dant8g3oe5ea%
}{Programme entry.}

\medskip\noindent
{\bf Abstract.}

A graph can be seen as the combinatorial
skeleton of a singular algebraic curve. On a
curve, we study line bundles and their spaces
of global sections via the Jacobian variety.
For singular curves, this Jacobian is not
compact, and finding a good compactification
is a deep problem in geometry. The
compactified Jacobian includes not only line
bundles on the curve but also on its partial
normalizations. In this talk, we explore how
this geometric idea can be imported into pure
combinatorics. We introduce a new invariant
for divisors on graphs, the balanced
compactified rank, defined by exploring all
curves that share the same combinatorial
skeleton and all line bundles in their
compactified Jacobians. We will see how this
new perspective, inspired by
compactification, allows us to refine
existing combinatorial invariants using
algebro-geometric intuition. We also prove
that this new rank satisfies a natural
Riemann-Roch theorem, and finish it with some
results on bridging the final gap between the
combinatorial and algebraic worlds,
suggesting that the geometry of
compactifications fully captures the
combinatorics of graphs. This is an ongoing
work with Margarida Melo.

\section{Codimension one distributions with
unstable tangent sheaf on $P^3$}
\label{section-codimension-one-distributions}
%% speaker: Pedro Pfarrius Barbassa, UNICAMP
%% speaker: - IMECC
%% place: LEGAL 2027, Sítio Assunção,
%% place: Teresópolis, RJ, Brazil

\medskip\noindent
{\bf Official presentation type.}

Contributed talk.

\medskip\noindent
{\bf Speaker.}

Pedro Pfarrius Barbassa, UNICAMP - IMECC.

\medskip\noindent
{\bf Schedule.}

Not listed on the official site as of
November 20, 2026.

\medskip\noindent
{\bf Official listing.}

\href{%
https://sites.google.com/view/%
geoalgcompluff/legal/%
legal-2027-resumos#h.dant8g3oe5ea%
}{Programme entry.}

\medskip\noindent
{\bf Abstract.}

We study codimension-one distributions on
$P^3$ with unstable tangent sheaf. The
instability naturally relates to an induced
one-dimensional subfoliation, allowing us to
the degree of this subfoliation to the order
of instability. We describe restrictions on
the possible instability, the resulting Chern
classes, and the geometry of the singular
locus. In the maximally unstable case, we
obtain a more precise description of the
tangent sheaf and of the associated singular
scheme.

\section{Multiple Tor modules and
Mayer-Vietoris spectral sequences}
\label{section-multiple-tor-modules-and}
%% speaker: Rafael Ferreira Holanda,
%% speaker: Universidade Federal de
%% speaker: Pernambuco
%% place: LEGAL 2027, Sítio Assunção,
%% place: Teresópolis, RJ, Brazil

\medskip\noindent
{\bf Official presentation type.}

Contributed talk.

\medskip\noindent
{\bf Speaker.}

Rafael Ferreira Holanda, Universidade Federal
de Pernambuco.

\medskip\noindent
{\bf Schedule.}

Not listed on the official site as of
November 20, 2026.

\medskip\noindent
{\bf Official listing.}

\href{%
https://sites.google.com/view/%
geoalgcompluff/legal/%
legal-2027-resumos#h.dant8g3oe5ea%
}{Programme entry.}

\medskip\noindent
{\bf Abstract.}

We extend some properties of a pair of ideals
described in terms of Tor modules to any
number of ideals. Those extensions require
the development of a homological theory for
spectral sequences arising from multiple
complexes. Out of this theory, two new
complexes associated with quotients by sums
and quotients by products of the given ideals
emerge. This investigation has consequences
in the study of extensions of the
Castelnuovo-Mumford regularity to a
multigraded setting, which is in part at the
origin of this work.

\section{Prym–Tyurin varieties from cyclic
covers}
\label{section-prymtyurin-varieties-from}
%% speaker: Victor Valdebenito Sepúlveda,
%% speaker: Universidad de la frontera
%% place: LEGAL 2027, Sítio Assunção,
%% place: Teresópolis, RJ, Brazil

\medskip\noindent
{\bf Official presentation type.}

Contributed talk.

\medskip\noindent
{\bf Speaker.}

Victor Valdebenito Sepúlveda, Universidad de
la frontera.

\medskip\noindent
{\bf Schedule.}

Not listed on the official site as of
November 20, 2026.

\medskip\noindent
{\bf Official listing.}

\href{%
https://sites.google.com/view/%
geoalgcompluff/legal/%
legal-2027-resumos#h.dant8g3oe5ea%
}{Programme entry.}

\medskip\noindent
{\bf Abstract.}

In this talk, a new construction of families
of Prym–Tyurin varieties arising from cyclic
covers of algebraic curves will be presented.
The construction provides explicit families
of principally polarized abelian varieties
and allows one to control the dimension of
their image in the corresponding moduli
space. Particular attention will be given to
the totally ramified case, where the geometry
can be described more explicitly. These
families are not obtained through Kanev’s
correspondence construction and, in the cases
considered, do not arise as Jacobians.

\section{Title to be announced: Wodson
Mendson}
\label{section-tba-mendson}
%% speaker: Wodson Mendson, UFF
%% place: LEGAL 2027, Sítio Assunção,
%% place: Teresópolis, RJ, Brazil

\medskip\noindent
{\bf Official presentation type.}

Contributed talk.

\medskip\noindent
{\bf Title.}

The official site lists the title as
``TBA'' as of November 20, 2026.

\medskip\noindent
{\bf Speaker.}

Wodson Mendson, UFF.

\medskip\noindent
{\bf Schedule.}

Not listed on the official site as of
November 20, 2026.

\medskip\noindent
{\bf Official listing.}

\href{%
https://sites.google.com/view/%
geoalgcompluff/legal/%
legal-2027-resumos#h.dant8g3oe5ea%
}{Programme entry.}

\medskip\noindent
{\bf Abstract.}

The official site lists the abstract
as ``TBA'' as of November 20, 2026.

\section{Rational Chow rings of Hurwitz
spaces with marked ramification}
\label{section-rational-chow-rings-of}
%% speaker: Zhengning Hu, UNICAMP
%% place: LEGAL 2027, Sítio Assunção,
%% place: Teresópolis, RJ, Brazil

\medskip\noindent
{\bf Official presentation type.}

Contributed talk.

\medskip\noindent
{\bf Speaker.}

Zhengning Hu, UNICAMP.

\medskip\noindent
{\bf Schedule.}

Not listed on the official site as of
November 20, 2026.

\medskip\noindent
{\bf Official listing.}

\href{%
https://sites.google.com/view/%
geoalgcompluff/legal/%
legal-2027-resumos#h.dant8g3oe5ea%
}{Programme entry.}

\medskip\noindent
{\bf Abstract.}

In this talk, we will discuss the computation
of the rational Chow rings of degree 3
Hurwitz spaces with marked ramification.
Aside from their intrinsic interest, products
of these spaces naturally arise as boundary
strata of the admissible covers
compactification of some Hurwitz spaces. For
each partition of 3, we focus on the moduli
space of degree 3 genus g covers over the
projective line which are simply branched
away from a special fiber with prescribed
ramification profile. By realizing each of
the marked Hurwitz spaces as an open substack
of a vector bundle over a well-understood
stack, we prove that they all have trivial
Chow rings. As a byproduct, the results imply
that the codimension 2 Chow group of the
admissible covers compactification of the
degree 3 genus g Hurwitz space is generated
by the classes of codimension 2 boundary
strata.

\section{Title to be announced: Adina
Veronica Remor}
\label{section-tba-poster-remor}
%% speaker: Adina Veronica Remor,
%% speaker: Universidade Federal do Paraná
%% place: LEGAL 2027, Sítio Assunção,
%% place: Teresópolis, RJ, Brazil

\medskip\noindent
{\bf Official presentation type.}

Poster.

\medskip\noindent
{\bf Title.}

The official site lists the title as
``TBA'' as of November 20, 2026.

\medskip\noindent
{\bf Speaker.}

Adina Veronica Remor, Universidade Federal do
Paraná.

\medskip\noindent
{\bf Schedule.}

Not listed on the official site as of
November 20, 2026.

\medskip\noindent
{\bf Official listing.}

\href{%
https://sites.google.com/view/%
geoalgcompluff/legal/%
legal-2027-resumos#h.f3btoqn8ykd7%
}{Programme entry.}

\medskip\noindent
{\bf Abstract.}

No abstract is available on the
official site as of November 20, 2026.

\section{Title to be announced: Ellen Peixoto
de Oliveira}
\label{section-tba-poster-oliveira}
%% speaker: Ellen Peixoto de Oliveira,
%% speaker: Unicamp
%% place: LEGAL 2027, Sítio Assunção,
%% place: Teresópolis, RJ, Brazil

\medskip\noindent
{\bf Official presentation type.}

Poster.

\medskip\noindent
{\bf Title.}

The official site lists the title as
``TBA'' as of November 20, 2026.

\medskip\noindent
{\bf Speaker.}

Ellen Peixoto de Oliveira, Unicamp.

\medskip\noindent
{\bf Schedule.}

Not listed on the official site as of
November 20, 2026.

\medskip\noindent
{\bf Official listing.}

\href{%
https://sites.google.com/view/%
geoalgcompluff/legal/%
legal-2027-resumos#h.f3btoqn8ykd7%
}{Programme entry.}

\medskip\noindent
{\bf Abstract.}

No abstract is available on the
official site as of November 20, 2026.

\section{Title to be announced: Joaquín
Esteban Lillo Godoy}
\label{section-tba-poster-godoy}
%% speaker: Joaquín Esteban Lillo Godoy,
%% speaker: Universidad de Chile
%% place: LEGAL 2027, Sítio Assunção,
%% place: Teresópolis, RJ, Brazil

\medskip\noindent
{\bf Official presentation type.}

Poster.

\medskip\noindent
{\bf Title.}

The official site lists the title as
``TBA'' as of November 20, 2026.

\medskip\noindent
{\bf Speaker.}

Joaquín Esteban Lillo Godoy, Universidad de
Chile.

\medskip\noindent
{\bf Schedule.}

Not listed on the official site as of
November 20, 2026.

\medskip\noindent
{\bf Official listing.}

\href{%
https://sites.google.com/view/%
geoalgcompluff/legal/%
legal-2027-resumos#h.f3btoqn8ykd7%
}{Programme entry.}

\medskip\noindent
{\bf Abstract.}

No abstract is available on the
official site as of November 20, 2026.

\section{Title to be announced: Larissa
Mariane dos Reis}
\label{section-tba-poster-reis}
%% speaker: Larissa Mariane dos Reis, UFMG
%% place: LEGAL 2027, Sítio Assunção,
%% place: Teresópolis, RJ, Brazil

\medskip\noindent
{\bf Official presentation type.}

Poster.

\medskip\noindent
{\bf Title.}

The official site lists the title as
``TBA'' as of November 20, 2026.

\medskip\noindent
{\bf Speaker.}

Larissa Mariane dos Reis, UFMG.

\medskip\noindent
{\bf Schedule.}

Not listed on the official site as of
November 20, 2026.

\medskip\noindent
{\bf Official listing.}

\href{%
https://sites.google.com/view/%
geoalgcompluff/legal/%
legal-2027-resumos#h.f3btoqn8ykd7%
}{Programme entry.}

\medskip\noindent
{\bf Abstract.}

No abstract is available on the
official site as of November 20, 2026.

\section{Title to be announced: Leonardo
Garcia de Castro}
\label{section-tba-poster-castro}
%% speaker: Leonardo Garcia de Castro,
%% speaker: UNICAMP
%% place: LEGAL 2027, Sítio Assunção,
%% place: Teresópolis, RJ, Brazil

\medskip\noindent
{\bf Official presentation type.}

Poster.

\medskip\noindent
{\bf Title.}

The official site lists the title as
``TBA'' as of November 20, 2026.

\medskip\noindent
{\bf Speaker.}

Leonardo Garcia de Castro, UNICAMP.

\medskip\noindent
{\bf Schedule.}

Not listed on the official site as of
November 20, 2026.

\medskip\noindent
{\bf Official listing.}

\href{%
https://sites.google.com/view/%
geoalgcompluff/legal/%
legal-2027-resumos#h.f3btoqn8ykd7%
}{Programme entry.}

\medskip\noindent
{\bf Abstract.}

No abstract is available on the
official site as of November 20, 2026.

\section{The Generalized Severi Variety:
Stratifying the Moduli of Linear Series on
$\mathbb{P}^1$ Using Value Semigroups}
\label{section-the-generalized-severi}
%% speaker: Lucas Henrique Martins da Silva,
%% speaker: UNICAMP
%% place: LEGAL 2027, Sítio Assunção,
%% place: Teresópolis, RJ, Brazil

\medskip\noindent
{\bf Official presentation type.}

Poster.

\medskip\noindent
{\bf Speaker.}

Lucas Henrique Martins da Silva, UNICAMP.

\medskip\noindent
{\bf Schedule.}

Not listed on the official site as of
November 20, 2026.

\medskip\noindent
{\bf Official listing.}

\href{%
https://sites.google.com/view/%
geoalgcompluff/legal/%
legal-2027-resumos#h.f3btoqn8ykd7%
}{Programme entry.}

\medskip\noindent
{\bf Abstract.}

No abstract is available on the
official site as of November 20, 2026.

\section{Title to be announced: Rodrigo
Contreras Martínez}
\label{section-tba-poster-martinez}
%% speaker: Rodrigo Contreras Martínez,
%% speaker: Universidad de Chile
%% place: LEGAL 2027, Sítio Assunção,
%% place: Teresópolis, RJ, Brazil

\medskip\noindent
{\bf Official presentation type.}

Poster.

\medskip\noindent
{\bf Title.}

The official site lists the title as
``TBA'' as of November 20, 2026.

\medskip\noindent
{\bf Speaker.}

Rodrigo Contreras Martínez, Universidad de
Chile.

\medskip\noindent
{\bf Schedule.}

Not listed on the official site as of
November 20, 2026.

\medskip\noindent
{\bf Official listing.}

\href{%
https://sites.google.com/view/%
geoalgcompluff/legal/%
legal-2027-resumos#h.f3btoqn8ykd7%
}{Programme entry.}

\medskip\noindent
{\bf Abstract.}

No abstract is available on the
official site as of November 20, 2026.

\section{Title to be announced: Sheucier
Alves de Medeiros}
\label{section-tba-poster-medeiros}
%% speaker: Sheucier Alves de Medeiros, UFMG
%% place: LEGAL 2027, Sítio Assunção,
%% place: Teresópolis, RJ, Brazil

\medskip\noindent
{\bf Official presentation type.}

Poster.

\medskip\noindent
{\bf Title.}

The official site lists the title as
``TBA'' as of November 20, 2026.

\medskip\noindent
{\bf Speaker.}

Sheucier Alves de Medeiros, UFMG.

\medskip\noindent
{\bf Schedule.}

Not listed on the official site as of
November 20, 2026.

\medskip\noindent
{\bf Official listing.}

\href{%
https://sites.google.com/view/%
geoalgcompluff/legal/%
legal-2027-resumos#h.f3btoqn8ykd7%
}{Programme entry.}

\medskip\noindent
{\bf Abstract.}

No abstract is available on the
official site as of November 20, 2026.

\section{Title to be announced: Valentina
Moreno Vega}
\label{section-tba-poster-vega}
%% speaker: Valentina Moreno Vega,
%% speaker: Universidad de Chile
%% place: LEGAL 2027, Sítio Assunção,
%% place: Teresópolis, RJ, Brazil

\medskip\noindent
{\bf Official presentation type.}

Poster.

\medskip\noindent
{\bf Title.}

The official site lists the title as
``TBA'' as of November 20, 2026.

\medskip\noindent
{\bf Speaker.}

Valentina Moreno Vega, Universidad de Chile.

\medskip\noindent
{\bf Schedule.}

Not listed on the official site as of
November 20, 2026.

\medskip\noindent
{\bf Official listing.}

\href{%
https://sites.google.com/view/%
geoalgcompluff/legal/%
legal-2027-resumos#h.f3btoqn8ykd7%
}{Programme entry.}

\medskip\noindent
{\bf Abstract.}

No abstract is available on the
official site as of November 20, 2026.

\section{Title to be announced: Ygor Allem
Rentroia}
\label{section-tba-poster-rentroia}
%% speaker: Ygor Allem Rentroia, UFF
%% place: LEGAL 2027, Sítio Assunção,
%% place: Teresópolis, RJ, Brazil

\medskip\noindent
{\bf Official presentation type.}

Poster.

\medskip\noindent
{\bf Title.}

The official site lists the title as
``TBA'' as of November 20, 2026.

\medskip\noindent
{\bf Speaker.}

Ygor Allem Rentroia, UFF.

\medskip\noindent
{\bf Schedule.}

Not listed on the official site as of
November 20, 2026.

\medskip\noindent
{\bf Official listing.}

\href{%
https://sites.google.com/view/%
geoalgcompluff/legal/%
legal-2027-resumos#h.f3btoqn8ykd7%
}{Programme entry.}

\medskip\noindent
{\bf Abstract.}

No abstract is available on the
official site as of November 20, 2026.

\bibliography{my}
\bibliographystyle{amsalpha}

\end{document}
